\documentclass[12pt,a4paper]{ctexart}

\usepackage[left=2.5cm,right=2.5cm,top=2.6cm,bottom=2.6cm]{geometry}
\usepackage{graphicx}
\usepackage{float}
\usepackage{xcolor}
\usepackage{listings}
\usepackage[hidelinks]{hyperref}
\usepackage{fancyhdr}
\usepackage{enumitem}

\pagestyle{fancy}
\fancyhf{}
\cfoot{\thepage}
\renewcommand{\headrulewidth}{0pt}

\setlength{\parindent}{2em}
\setlength{\parskip}{0.35em}
\linespread{1.25}

\lstset{
  basicstyle=\ttfamily\small,
  breaklines=true,
  columns=fullflexible,
  keepspaces=true,
  frame=single,
  numbers=none,
  showstringspaces=false,
  tabsize=4
}

% listings does not provide JSON by default; define it explicitly so the
% JSON example remains portable across TeX Live installations.
\lstdefinelanguage{json}{}

\newcommand{\screenshot}[2]{
\begin{figure}[H]
  \centering
  \includegraphics[width=0.94\textwidth]{report_assets/practice/#1}
  \caption{#2}
\end{figure}
}

\begin{document}

\begin{titlepage}
\thispagestyle{empty}
\begin{center}
\vspace*{3.2cm}

{\zihao{1}\bfseries 系统开发工具基础第四周实验报告}

\vfill

{\zihao{3}
姓名：刘静雅\\[0.8cm]
学号：24160001055\\[0.8cm]
日期：2026年9月14日
}

\vspace*{4.2cm}
\end{center}
\end{titlepage}

\clearpage
\thispagestyle{plain}
\tableofcontents
\clearpage

\section{实验环境与目录结构}

本周实验在 WSL 的 Bash 环境中完成，仓库目录为
\verb|~/sys-dev-tools-lab/week04|。实验主要使用 Python、pytest、ruff、
Make、curl、jq、Git 等工具，完成代码质量检查、自动化构建、本地 API
数据处理与软件包交付等任务。

本周实验目录结构如下：

\begin{lstlisting}[language=bash]
week04/
|-- q13/
|-- q14/
|-- q15/
|-- q16/
|-- report_assets/
|   `-- practice/
`-- report.tex
\end{lstlisting}

\section{课上实验}

\subsection{第13题： 建立可执行的本地质量门禁}

\subsubsection*{题目原文}

复用第10题完成的 greetlab 包，在 q13 中建立一个轻量本地质量检查。

\begin{itemize}
  \item 复制 q10 为 q13；在 pyproject.toml 中加入 ruff 和 pytest 的最小配置。
  \item 补充一个正常姓名测试和一个空白姓名测试，每个测试包含有效断言。
  \item 运行 ruff format、ruff check 和 pytest，修复全部问题，不得全局忽略规则。
  \item 编写 check.sh，依次执行 ruff format --check、ruff check 和 pytest；运行并保证返回 0。
\end{itemize}

\subsubsection*{实验步骤}

首先复用前一周完成的 greetlab 项目，将 q10 复制为 q13。在
\verb|pyproject.toml| 中补充 pytest 与 ruff 的最小配置，使测试框架能够
定位 \verb|src| 与 \verb|tests| 目录，同时为代码格式检查和静态检查提供统一规则。

随后补充正常姓名和空白姓名两个测试用例。正常姓名测试验证程序能够输出正确
问候语；空白姓名测试验证程序能够拒绝无效输入。完成代码修改后，分别运行
\verb|pytest|、\verb|ruff format --check .| 和 \verb|ruff check|。

最后编写 \verb|check.sh|，将格式检查、静态检查和单元测试串联为一个统一的
本地质量门禁入口。脚本的核心形式如下：

\begin{lstlisting}[language=bash]
#!/usr/bin/env bash
set -euo pipefail

ruff format --check .
ruff check .
pytest
\end{lstlisting}

\subsubsection*{实验结果与分析}

\screenshot{practice13_code.png}{第13题 pyproject.toml 中的 pytest 与 ruff 配置}

实验完成后，\verb|pyproject.toml| 中已经配置了 pytest 的
\verb|pythonpath| 与 \verb|testpaths|，并设置了 ruff 的目标 Python 版本、
行长度和检查规则。这些配置使质量检查可以在项目目录中稳定执行。

\screenshot{practice13_quality_pass.png}{第13题 pytest、ruff format 与 ruff check 通过}

执行检查命令后，pytest 共收集两个测试项目，最终结果为 ``2 passed''。随后
\begin{lstlisting}[language=bash]
ruff format --check .
\end{lstlisting}
显示全部文件格式正确，\verb|ruff check| 显示
``All checks passed!''，说明测试、格式与静态检查均通过。

\screenshot{practice13_check_sh.png}{第13题 check.sh 执行结果与退出码}

执行 \verb|check.sh| 后，格式检查、静态检查和测试均成功完成，最后
\verb|echo $?| 输出为 0。返回值 0 表明该质量门禁脚本执行成功，可在提交前
作为统一检查入口使用。

\subsection{第14题： 让 Make 只重建真正受影响的产物}

\subsubsection*{题目原文}

在 q14 中按给定内容创建三个源文件，只需完成 Makefile。

给定内容如下：

\begin{lstlisting}[language=bash]
# data.csv
name,value
alpha,2
beta,3
gamma,5

# stats.py
import csv
with open("data.csv") as f:
    total = sum(int(r["value"]) for r in csv.DictReader(f))
open("stats.txt", "w").write(str(total))

# build_report.py
text = open("report.md").read()
total = open("stats.txt").read()
open("report.txt", "w").write(f"{text}\nTotal: {total}\n")

# report.md
# Data Report
\end{lstlisting}

\begin{itemize}
  \item 创建给定的 data.csv、stats.py、build\_report.py 和 report.md。
  \item 编写 Makefile，包含 all、stats.txt、report.txt 和 clean；准确列出数据文件、脚本和中间产物依赖，并把 all、clean 声明为 .PHONY。
  \item 验证首次 make 生成两个产物；第二次无改动时不执行配方。
  \item touch data.csv 后再次 make，确认 stats.txt 和 report.txt 重新生成；clean 只删除生成物。
\end{itemize}

\subsubsection*{实验步骤}

本题先创建数据文件 \verb|data.csv|、统计脚本 \verb|stats.py|、报告生成脚本
\verb|build_report.py| 以及报告标题文件 \verb|report.md|。其中，
\verb|stats.py| 用于计算 CSV 文件中数值列的总和，\verb|build_report.py| 用于
读取统计结果并生成最终报告。

随后编写 Makefile，明确目标文件之间的依赖关系。其核心内容如下：

\begin{lstlisting}[language=make]
.PHONY: all clean

all: report.txt

stats.txt: data.csv stats.py
	python3 stats.py

report.txt: stats.txt report.md build_report.py
	python3 build_report.py

clean:
	rm -f stats.txt report.txt
\end{lstlisting}

通过首次执行 \verb|make|、再次执行 \verb|make|、执行
\verb|touch data.csv| 后重新执行 \verb|make|，可以验证 Make 的增量构建机制。
最后使用 \verb|make clean| 删除构建产生的中间文件与结果文件。

\subsubsection*{实验结果与分析}

\screenshot{practice14_first_make.png}{第14题首次 make 构建及生成结果}

首次执行 \verb|make| 时，终端依次运行 \verb|python3 stats.py| 和
\verb|python3 build_report.py|。生成的 \verb|stats.txt| 内容为 10，
\verb|report.txt| 中包含报告标题和 ``Total: 10''，说明数据统计与报告生成均成功完成。

\screenshot{practice14_no_rebuild.png}{第14题第二次 make 未重复构建}

第二次直接运行 \verb|make| 时，终端显示
``Nothing to be done for `all'.''。这表示源文件没有发生变化，
目标文件仍然是最新状态，因此 Make 没有重复执行构建命令。

\screenshot{practice14_rebuild_after_touch.png}{第14题更新 data.csv 后重新构建}

执行 \verb|touch data.csv| 后再次运行 \verb|make|，终端重新执行
\verb|stats.py| 与 \verb|build_report.py|。这是因为数据文件的时间戳更新后，
\verb|stats.txt| 以及依赖它的 \verb|report.txt| 都需要重新生成。

\screenshot{practice14_clean.png}{第14题 make clean 清理结果}

执行 \verb|make clean| 后，终端删除 \verb|stats.txt| 和 \verb|report.txt|。
随后列出的目录文件中仅保留 Makefile、源数据文件、Python 脚本和 Markdown 文件，
说明 clean 目标只清理生成物，没有删除源文件。

\subsection{第15题： 把本地 API 数据转换为可读报告}

\subsubsection*{题目原文}

在 q15 中创建给定 JSON 和本地 HTTP 服务，再用 curl 与 jq 生成简短 Markdown 报告。

给定 JSON 如下：

\begin{lstlisting}
[
  {"name":"alpha","status":"active","downloads":120,"version":"1.2.0"},
  {"name":"beta","status":"inactive","downloads":900,"version":"2.0.0"},
  {"name":"gamma","status":"active","downloads":450,"version":"1.5.1"},
  {"name":"delta","status":"active","downloads":450,"version":"0.9.0"},
  {"name":"epsilon","status":"active","downloads":80,"version":"3.1.0"}
]
\end{lstlisting}

\begin{itemize}
  \item 把给定 JSON 保存为 packages.json，并在 q15 目录运行 python -m http.server 8000。
  \item 使用 curl -fsS 获取 http://127.0.0.1:8000/packages.json。
  \item 使用 jq 筛选 status 为 active 且 downloads 不少于 100 的记录，并按 downloads 降序、name 升序排列。
  \item 编写 api\_report.sh，把结果生成 summary.md，包含标题和 name/version/downloads 三列表格。
\end{itemize}

\subsubsection*{实验步骤}

首先将题目给定的 JSON 数据保存为 \verb|packages.json|，并在 q15 目录启动
本地 HTTP 服务。随后使用 \verb|curl -fsS| 获取本地服务器返回的 JSON 数据，
再使用 jq 完成筛选、排序和格式化。

脚本 \verb|api_report.sh| 先输出 Markdown 标题与表头，再对 JSON 数组执行
筛选和排序，最终生成 \verb|summary.md|。筛选逻辑的关键部分如下：

\begin{lstlisting}[language=bash]
map(select(.status == "active" and .downloads >= 100))
| sort_by(-.downloads, .name)
| .[]
| "| \(.name) | \(.version) | \(.downloads) |"
\end{lstlisting}

筛选后排除 inactive 的 beta 和 downloads 小于 100 的 epsilon。delta 和
gamma 的下载量同为 450，因此按照名称升序排列，最终顺序为 delta、gamma、alpha。

\subsubsection*{实验结果与分析}

\screenshot{practice15_curl_jq.png}{第15题 curl 获取本地 JSON 数据}

执行 \verb|curl -fsS http://127.0.0.1:8000/packages.json| 后，命令能够成功访问
本地 HTTP 服务，jq 将 JSON 数据格式化输出。该结果说明本地 API 服务和数据文件
均可正常访问。

\screenshot{practice15_script.png}{第15题 api\_report.sh 脚本内容}

脚本使用 \verb|set -euo pipefail| 保证命令执行出现错误时立即停止。
脚本通过 curl 获取 JSON，再使用 jq 筛选 active 状态且 downloads 不少于 100 的
记录，并将结果写入 \verb|summary.md|。

\screenshot{practice15_summary.png}{第15题 summary.md 报告结果}

生成的 Markdown 报告包含 name、version、downloads 三列，最终保留
delta、gamma 和 alpha 三条记录。结果中的排序符合下载量降序、名称升序的要求。

\subsection{第16题： 修复并交付一个陌生的小型工具仓库}

\subsubsection*{题目原文}

复用第13题的项目，在 q16 中完成一次 15 分钟综合检查。

\begin{itemize}
  \item 复制 q13 为 q16 并初始化/继续使用 Git。
  \item 把问候语实现临时改成始终返回字面量 Hello, name!，运行 check.sh 确认测试失败；随后定位并修复。
  \item 编写最小 Makefile，包含 check、build 和 clean；check 调用 check.sh，build 生成 wheel。
  \item 运行 make check 和 make build，计算 wheel 的 SHA-256，并把最终改动提交为一个内容聚焦的 Git 提交。
\end{itemize}

\subsubsection*{实验步骤}

本题首先将 q13 复制为 q16，并继续使用仓库根目录已有的 Git 管理。随后按题目要求，
将问候语实现临时改为固定输出 \verb|Hello, name!|，运行 \verb|check.sh| 以复现
测试失败。

在确认失败后，将输出逻辑修复为根据参数中的姓名输出问候语：

\begin{lstlisting}[language=python]
print(f"Hello, {name}!")
\end{lstlisting}

修复完成后重新执行 \verb|check.sh|，确认代码质量检查和两个测试用例均通过。

之后编写最小 Makefile，将检查、构建与清理任务封装为 \verb|make| 目标：

\begin{lstlisting}[language=make]
.PHONY: check build clean

check:
	./check.sh

build:
	python -m build

clean:
	rm -rf build dist src/*.egg-info
\end{lstlisting}

最后运行 \verb|make build| 构建源码包和 wheel 包，并通过
\verb|sha256sum dist/*.whl| 计算 wheel 文件的 SHA-256 校验值。

\subsubsection*{实验结果与分析}

\screenshot{practice16_intentional_fail.png}{第16题故意修改问候语后测试失败}

将问候语临时修改为固定文本 \verb|Hello, name!| 后，重新运行测试，正常姓名测试失败。
测试期望输出为 \verb|Hello, Ada!|，但实际输出为 \verb|Hello, name!|，因此触发
断言错误。该结果说明已有测试能够有效发现功能回归。

\screenshot{practice16_check_pass.png}{第16题修复后 check.sh 通过}

恢复为基于姓名变量输出问候语后，执行 \verb|check.sh| 的结果显示两个测试全部通过。
终端输出 \verb|exit_code=0|，说明格式检查、静态检查和测试均执行成功。

\screenshot{practice16_makefile.png}{第16题 Makefile 内容}

Makefile 中定义了 \verb|check|、\verb|build| 和 \verb|clean| 三个目标。
其中 check 调用 \verb|check.sh|，build 调用 \verb|python -m build|，clean 用于删除
构建产生的 build、dist 和 egg-info 文件夹。

\screenshot{practice16_make_build.png}{第16题 make build 构建过程}

执行 \verb|make build| 后，命令成功调用 \verb|python -m build|，并完成源码包构建流程。
构建过程中出现 README 文件缺失的警告，但没有导致构建中断，后续仍成功生成了发布产物。

\screenshot{practice16_build_sha256.png}{第16题 wheel 构建完成并计算 SHA-256}

构建过程最终成功生成 \verb|.tar.gz| 与 \verb|.whl| 文件。随后使用
\verb|sha256sum dist/*.whl| 输出 wheel 文件的 SHA-256 校验值。该值可用于验证
构建产物的完整性，便于后续提交与交付时进行核验。

\section{课后练习}

本部分按照实际完成的十个练习记录实验过程。每题先列出题目要求，再说明实验方法，最后结合运行结果进行分析。练习均在独立的质量实践项目中完成，相关配置位于仓库的 \verb|week04/practice/quality-lab| 目录。

\subsection{代码格式化、静态检查与 AI 修复}
\subsubsection*{题目原文}
为正在进行的项目配置格式化器、linter 和 AI agent。自动格式化应当能够处理格式错误；对于 linter 错误，要求 AI agent 运行 linter 并根据结果迭代修复。修复完成后需要人工检查改动，确认代码功能没有被破坏。

\subsubsection*{实验步骤与分析}
在质量实践项目中安装 Ruff，先执行以下检查命令并记录原始结果：
\begin{lstlisting}[language=bash]
ruff format --check .
ruff check .
\end{lstlisting}
随后把检查信息和修复目标交给 AI agent，要求它运行相同的检查命令，只修改必要文件。修复结束后再次执行检查，并用 \verb|git diff| 人工核对改动范围。自动格式化适合处理缩进、换行等机械性问题，而 linter 提示仍需结合代码语义确认。

\subsubsection*{实验结果与解释}
\screenshot{practice01_ruff_before_ai.png}{课后练习一：AI 修复前的 Ruff 检查}
\screenshot{practice01_ai_fix_and_review.png}{课后练习一：AI 修复后的检查与人工复核}
初始检查能够定位格式和静态检查问题；修复后再次运行 Ruff，检查通过，差异内容集中在格式调整和 linter 指出的代码位置。这说明 AI agent 能够参与修复循环，但最终结果仍需人工复核。

\subsection{提交前自动检查}
\subsubsection*{题目原文}
为项目配置 pre-commit 钩子，使提交前自动运行格式化检查、linter 和测试。

\subsubsection*{实验步骤与分析}
在项目根目录创建 \verb|.pre-commit-config.yaml|，将格式检查、静态检查和测试配置为本地钩子。安装 Git 钩子并检查全部文件：
\begin{lstlisting}[language=bash]
pre-commit install
pre-commit run --all-files
\end{lstlisting}
这样可以把原本需要手动执行的质量检查接入提交流程。

\subsubsection*{实验结果与解释}
\screenshot{practice02_precommit_pass.png}{课后练习二：pre-commit 全部检查通过}
全量检查完成后，格式检查、静态检查和测试均通过。结果表明，pre-commit 能够在提交前统一执行质量检查，减少遗漏手工检查的可能。

\subsection{pytest 单元测试}
\subsubsection*{题目原文}
学习一种熟悉语言的测试库，并为正在进行的项目编写一个单元测试，运行测试并根据结果判断程序行为是否符合预期。

\subsubsection*{实验步骤与分析}
在现有 Python 项目中补充 pytest 测试，分别覆盖正常输入、边界输入和无效输入，并通过断言验证函数返回值或命令行输出。执行 \verb|pytest -v|，确认测试文件被正确收集，并根据失败信息修正实现或测试条件。

\subsubsection*{实验结果与解释}
\screenshot{practice03_unit_test.png}{课后练习三：pytest 单元测试结果}
测试命令返回成功，说明测试文件被正确识别，断言全部成立。正常输入和异常输入分别验证了主要功能路径与错误处理路径。

\subsection{覆盖率统计与 HTML 报告}
\subsubsection*{题目原文}
运行代码覆盖率工具并生成 HTML 格式报告，观察哪些代码行已经被测试覆盖，找出尚未覆盖的代码位置。

\subsubsection*{实验步骤与分析}
使用以下命令同时生成终端统计与 HTML 报告：
\begin{lstlisting}[language=bash]
pytest --cov=greetlab \
  --cov-report=term-missing \
  --cov-report=html
\end{lstlisting}
先查看终端中的覆盖率和缺失行，再打开 HTML 报告定位具体文件和代码行，为后续补充测试提供依据。

\subsubsection*{实验结果与解释}
\screenshot{practice04_coverage_initial.png}{课后练习四：初始覆盖率统计}
\screenshot{practice04_html_initial.png}{课后练习四：HTML 覆盖率报告}
初始覆盖率为 94\%，HTML 报告标出了尚未执行的代码行。覆盖率不能直接证明程序没有错误，但能够明确指出测试没有涉及的路径。

\subsection{人工补充测试与覆盖率提升}
\subsubsection*{题目原文}
根据覆盖率报告中未覆盖的代码行，手动编写测试以提高覆盖率，并重新运行测试确认新增测试有效。

\subsubsection*{实验步骤与分析}
根据 HTML 报告定位未覆盖分支，分析这些分支对应的输入条件，然后在测试文件中补充针对性的边界测试或异常测试。重新运行带覆盖率的 pytest，并检查新增测试是否包含有效断言，而不是只为了执行代码。

\subsubsection*{实验结果与解释}
\screenshot{practice05_manual_coverage.png}{课后练习五：补充测试后的覆盖率结果}
补充测试后，原先未执行的分支被实际运行，测试仍然全部通过，覆盖率得到提高。这说明提高覆盖率的关键是覆盖新的行为路径，而不是单纯增加测试数量。

\subsection{AI agent 辅助测试改进}
\subsubsection*{题目原文}
使用 AI agent 提高项目覆盖率。要求 agent 以带覆盖率的方式运行测试，生成逐行覆盖率结果，并根据结果补充测试；同时人工判断 AI 生成的测试是否合理。

\subsubsection*{实验步骤与分析}
将带有逐行覆盖率的测试结果提供给 AI agent，要求它先运行以下命令，再根据缺失行提出并实现测试：
\begin{lstlisting}[language=bash]
pytest --cov=greetlab \
  --cov-report=term-missing
\end{lstlisting}
生成测试后重新运行覆盖率命令，并检查测试是否验证真实行为，是否包含清晰断言。

\subsubsection*{实验结果与解释}
\screenshot{practice06_ai_coverage.png}{课后练习六：AI agent 辅助补充测试后的覆盖率}
AI agent 根据覆盖率结果补充了测试，重新运行后测试通过且覆盖率得到改善。人工检查确认新增测试具有实际输入和断言，说明 AI 可以提高编写效率，但不能替代测试设计与复核。

\subsection{GitHub Actions 持续集成}
\subsubsection*{题目原文}
为项目设置持续集成，使其在每次 push 时运行格式化检查、lint 和测试。

\subsubsection*{实验步骤与分析}
在 \verb|.github/workflows/week04-quality.yml| 中配置 push 触发的工作流，指定 Python 版本，安装项目依赖，并依次执行格式检查、Ruff 检查和带覆盖率的 pytest。提交并推送后，在 GitHub Actions 页面查看运行状态。

\subsubsection*{实验结果与解释}
\screenshot{practice07_ci_pass.png}{课后练习七：GitHub Actions 检查通过}
推送后，工作流完成环境准备、依赖安装、格式检查、lint 和测试，最终运行成功。这表明项目在干净的远程环境中也能通过质量检查。

\subsection{持续集成失败与恢复}
\subsubsection*{题目原文}
故意破坏代码，例如引入一条 linter 违规，确认持续集成能够捕获问题；随后修复错误并再次推送，确认 CI 恢复成功。

\subsubsection*{实验步骤与分析}
先在代码中加入一条可检测的 lint 违规，提交并推送，观察 CI 是否在质量检查阶段失败。随后撤销该错误，再次运行本地检查并推送修复结果，比较两次工作流状态。

\subsubsection*{实验结果与解释}
\screenshot{practice08_ci_failure.png}{课后练习八：故意引入问题后的 CI 失败}
\screenshot{practice08_ci_recovered.png}{课后练习八：修复问题后的 CI 恢复}
引入 lint 违规后，CI 检查失败；修复并重新推送后，工作流恢复通过。前后结果证明工作流实际执行了质量检查，而不是只创建了配置文件。

\subsection{正则检查与 Markdown 替换}
\subsubsection*{题目原文}
写一个正则表达式，使用 grep 查找代码中 \verb|subprocess.Popen(..., shell=True)| 的出现位置，再破坏正则模式，比较 grep 与 Semgrep 的匹配能力。同时在 Markdown 文档中使用正则查找替换，将项目符号 \verb|-| 替换为 \verb|*|，但不能误改其他用途的短横线。

\subsubsection*{实验步骤与分析}
先在测试文件中写入单行和跨行的危险调用，使用 \verb|grep -nE| 检查简单模式，再使用 Semgrep 规则进行结构化匹配。随后使用 \verb|^(\s*)-\s+| 查找行首项目符号，并将捕获的缩进保留后替换为星号。该限制条件可以避免误改普通文本中的连字符。

\subsubsection*{实验结果与解释}
\screenshot{practice09_grep_semgrep.png}{课后练习九：grep 与 Semgrep 检查危险调用}
\screenshot{practice09_markdown_replace.png}{课后练习九：Markdown 项目符号正则替换}
grep 能匹配同一行的危险调用，但对跨行写法不够稳定；Semgrep 根据语法结构识别危险调用。Markdown 替换只处理行首项目符号，保留了普通文本中的短横线。

\subsection{JSON 字段提取与解析器}
\subsubsection*{题目原文}
写一个正则表达式，从形如 \verb|{"name": "Alyssa P. Hacker", "college": "MIT"}| 的 JSON 结构中捕获 name。处理 name 中包含转义双引号的情况，并使用 JSON 解析器从标准输入读取数据并输出 name。

\subsubsection*{实验步骤与分析}
首先使用带有非贪婪和转义处理的正则提取 name 字段：
\begin{lstlisting}
"name"\s*:\s*"((?:\\.|[^"\\])*)"
\end{lstlisting}
该模式允许字段值中出现经过 JSON 转义的双引号。随后使用 Python 的 \verb|json| 模块读取标准输入，并通过 \verb|json.load| 访问 name 字段。解析器能够遵守完整 JSON 语法，适合实际实现。

\subsubsection*{实验结果与解释}
\screenshot{practice10_regex.png}{课后练习十：正则提取 JSON 中的 name 字段}
\screenshot{practice10_json_parser.png}{课后练习十：使用 Python JSON 解析器输出 name}
正则方法能够提取普通姓名并处理转义双引号；JSON 解析器输出了相同字段值。比较两种方法可以看出，正则适合演示字符串模式匹配，而实际处理结构化 JSON 时应优先使用专用解析器。
\section{实验总结与感悟}

本周实验围绕代码质量、增量构建、数据处理、软件交付和自动化协作展开。第13题通过 pytest、
ruff 和 check.sh 建立了可重复执行的本地质量门禁，使代码格式、静态检查和功能测试
能够由统一命令完成。第14题通过 Makefile 声明文件之间的依赖关系，验证了 Make
能够根据源文件变化只重建必要产物，从而提高构建效率。

第15题将 JSON 文件通过本地 HTTP 服务提供，再结合 curl 与 jq 完成数据获取、
筛选、排序和 Markdown 报告生成，体现了命令行工具在数据处理任务中的组合能力。
第16题则将错误复现、测试失败、代码修复、质量验证、构建 wheel 和计算 SHA-256
校验值串联起来，形成了较完整的软件交付流程。

课后练习进一步补充了 pre-commit、覆盖率、持续集成、Semgrep 和 JSON 解析等方法。
这些练习使我认识到，质量控制应当贯穿开发、测试、提交和交付的全过程；AI agent
可以参与格式修复和测试设计，但其结果仍需通过命令执行、覆盖率数据和人工审查进行
验证。代码不仅应当能够运行，还应当具备明确的测试依据、可重复的构建过程和可验证
的交付产物。

\end{document}

